% Notation and Conventions
% Converted on 2026-10-08 from docs/lswt/00-foundations/notation-and-conventions.md (status: in-review, source-section: Summary of Notation and Symbols).
% Review status: in-review; awaiting the user's physics and mathematics review.

\hypertarget{sec-lswt-notation-and-conventions}{%
\section{Notation and Conventions}\label{sec-lswt-notation-and-conventions}}

The notation below distinguishes spin sites, magnetic unit cells, interaction matrices, and bosonic operators throughout the LSWT derivation. Each symbol has a fixed meaning within its stated domain. Physical definitions and derivations are developed in the linked notes; the source-to-current notation correspondence is listed under \protect\hyperlink{legacy-source-mapping}{Legacy Source Mapping}.

\subsection{Typography}

\begin{longtable}[]{@{}
  >{\raggedright\arraybackslash}p{(\columnwidth - 4\tabcolsep) * \real{0.3333}}
  >{\raggedright\arraybackslash}p{(\columnwidth - 4\tabcolsep) * \real{0.3333}}
  >{\raggedright\arraybackslash}p{(\columnwidth - 4\tabcolsep) * \real{0.3333}}@{}}
\toprule\noalign{}
\begin{minipage}[b]{\linewidth}\raggedright
Object
\end{minipage} & \begin{minipage}[b]{\linewidth}\raggedright
Form
\end{minipage} & \begin{minipage}[b]{\linewidth}\raggedright
Examples
\end{minipage} \\
\midrule\noalign{}
\endhead
\bottomrule\noalign{}
\endlastfoot
Scalar & Italic & \(S_I\), \(E_{\mathrm{cl}}\), \(\varepsilon_{n\mathbf{k}}\) \\
Vector & Bold & \(\mathbf{r}_I\), \(\mathbf{k}\), \(\hat{\mathbf{S}}_I\) \\
Latin-letter matrix & Sans serif & \(\mathsf{J}_{\ell}\), \(\mathsf{D}_I\), \(\mathsf{R}_I\), \(\mathsf{H}_{\mathbf{k}}\) \\
Quantum operator & Hat & \(\hat H\), \(\hat a_{i\mu}\), \(\hat\Psi_{\mathbf{k}}\) \\
Set & Calligraphic & \(\mathcal{L}\) \\
Identity matrix & Sans serif & \(\mathsf{I}_N\) \\
\end{longtable}

Hats denote quantum operators. Classical spin vectors and unit directions are written without hats, as \(\mathbf S_I\) and \(\mathbf n_I\). The Hamiltonian operator \(\hat H\) is therefore distinct from its quadratic coefficient matrix \(\mathsf H_{\mathbf k}\).

Conventional Greek matrix symbols, such as \(\Sigma_3\), are exceptions to the sans-serif rule for Latin letters. Their dimensions and roles are stated where they are introduced. Transposition, complex conjugation, and Hermitian conjugation are denoted by superscripts \(\mathsf T\), \(*\), and \(\dagger\), respectively. Exponentials are written as \(\exp(\cdots)\), and the imaginary unit is \(\mathrm i\).

A base letter may be shared by different object classes when typography makes the distinction explicit. Examples include the temperature \(T\) and transformation matrix \(\mathsf T_{\mathbf k}\), or the Dzyaloshinskii--Moriya vector \(\mathbf D_\ell\) and on-site anisotropy matrix \(\mathsf D_I\). The same matrix symbol is not assigned two different meanings.

\subsection{Indices and System Size}

\begin{longtable}[]{@{}
  >{\raggedright\arraybackslash}p{(\columnwidth - 4\tabcolsep) * \real{0.3333}}
  >{\raggedright\arraybackslash}p{(\columnwidth - 4\tabcolsep) * \real{0.3333}}
  >{\raggedright\arraybackslash}p{(\columnwidth - 4\tabcolsep) * \real{0.3333}}@{}}
\toprule\noalign{}
\begin{minipage}[b]{\linewidth}\raggedright
Index
\end{minipage} & \begin{minipage}[b]{\linewidth}\raggedright
Meaning
\end{minipage} & \begin{minipage}[b]{\linewidth}\raggedright
Domain
\end{minipage} \\
\midrule\noalign{}
\endhead
\bottomrule\noalign{}
\endlastfoot
\(i,j\) & Magnetic unit-cell index & Real space \\
\(\mu,\nu\) & Magnetic sublattice index & One magnetic unit cell \\
\(I,J\) & Physical spin-site index & \(I=(i,\mu)\), \(J=(j,\nu)\) \\
\(\ell\) & Representative physical link & \(\ell=(I,J)\in\mathcal L\) \\
\(\alpha,\beta,\gamma\) & Laboratory-frame Cartesian component & \(x,y,z\) \\
\(n,m\) & Positive-energy magnon-band index & Momentum space \\
\end{longtable}

The indices \(n,m\) are reserved for magnon bands, so a sublattice index \(\mu\) remains distinct from a band index \(n\). Local circular components are labeled \(+,-,0\); these labels do not replace the Cartesian meanings of \(\alpha,\beta,\gamma\). The local basis is defined in Classical Order and Local Frame (Section~\ref{sec-lswt-classical-order-and-local-frame}).

The inverse temperature is

\[
\beta\equiv(k_{\mathrm B}T)^{-1}.
\]

When a Cartesian index \(\beta\) also appears in the same expression or immediately adjacent discussion, the inverse temperature is written as \(\beta_T\). This exception distinguishes an index from a scalar without reassigning \(i,j\) as Cartesian indices. The thermodynamic entropy is denoted by \(\mathcal S\) to distinguish it from the spin length \(S_I\); its definition belongs to \lswtnoteref{LSWT observables}{Thermodynamics}.

System sizes are denoted by

\begin{longtable}[]{@{}
  >{\raggedright\arraybackslash}p{(\columnwidth - 2\tabcolsep) * \real{0.5000}}
  >{\raggedright\arraybackslash}p{(\columnwidth - 2\tabcolsep) * \real{0.5000}}@{}}
\toprule\noalign{}
\begin{minipage}[b]{\linewidth}\raggedright
Symbol
\end{minipage} & \begin{minipage}[b]{\linewidth}\raggedright
Meaning
\end{minipage} \\
\midrule\noalign{}
\endhead
\bottomrule\noalign{}
\endlastfoot
\(N_{\mathrm{uc}}\) & Number of magnetic unit cells \\
\(N_{\mathrm{sub}}\) & Number of magnetic sublattices per magnetic unit cell \\
\(N_{\mathrm{site}}\) & Total number of physical spin sites \\
\(N_k\) & Number of discrete momentum points \\
\end{longtable}

so that

\[
N_{\mathrm{site}}
=N_{\mathrm{uc}}N_{\mathrm{sub}},
\qquad
\dim\mathsf H_{\mathbf k}=2N_{\mathrm{sub}}.
\]

The source symbols \(L\), \(N\), and \(m_s\) are replaced by these quantities because their meanings vary across the source sections.

\subsection{Sites and Geometry}

The basis position of sublattice \(\mu\) within a magnetic unit cell is \(\boldsymbol\delta_\mu\). A physical site and its position are written as

\[
I=(i,\mu),
\qquad
\mathbf r_I=\mathbf R_i+\boldsymbol\delta_\mu,
\]

where \(\mathbf R_i\) is a magnetic Bravais-lattice vector. Equal classical spin angles alone do not identify two sites as the same magnetic sublattice; the magnetic translation pattern and local environment also enter the sublattice assignment.

For a two-dimensional lattice, the primitive vectors \(\mathbf a_1,\mathbf a_2\) and reciprocal vectors \(\mathbf b_1,\mathbf b_2\) satisfy

\[
\mathbf a_r\cdot\mathbf b_s=2\pi\delta_{rs},
\qquad r,s\in\{1,2\}.
\]

Real-space vectors share one Cartesian coordinate system, and momenta are expressed in the corresponding reciprocal Cartesian coordinates.

\subsection{Links and Bond Orientation}

The set \(\mathcal L\) contains each physical inter-site bond once. The representative orientation \(\ell=(I,J)\) fixes the bond displacement and the ordering of the exchange matrix. The reverse orientation \(\bar\ell=(J,I)\) is not separately included in \(\mathcal L\).

The displacement from the first endpoint to the second is

\[
\boldsymbol\Delta_\ell
=\mathbf r_J-\mathbf r_I
=(\mathbf R_j-\mathbf R_i)
 +(\boldsymbol\delta_\nu-\boldsymbol\delta_\mu),
\qquad
\ell=((i,\mu),(j,\nu)).
\]

The laboratory-frame exchange matrix is \(\mathsf J_\ell\), with Cartesian components \(J_\ell^{\alpha\beta}\). Reversing the bond orientation gives

\[
\mathsf J_{\bar\ell}=\mathsf J_\ell^{\mathsf T},
\qquad
J_{JI}^{\beta\alpha}=J_{IJ}^{\alpha\beta}.
\]

This relation uses the real Cartesian exchange matrices of the bilinear model. Its expression in a complex circular basis depends on the component convention, which is not fixed here. The relation between the link sum and an ordered-pair sum, together with the separate treatment of on-site interactions, is given in Bilinear Spin Hamiltonian (Section~\ref{sec-lswt-bilinear-spin-hamiltonian}).

\subsection{Spins and Local Frames}

\begin{longtable}[]{@{}
  >{\raggedright\arraybackslash}p{(\columnwidth - 2\tabcolsep) * \real{0.5000}}
  >{\raggedright\arraybackslash}p{(\columnwidth - 2\tabcolsep) * \real{0.5000}}@{}}
\toprule\noalign{}
\begin{minipage}[b]{\linewidth}\raggedright
Symbol
\end{minipage} & \begin{minipage}[b]{\linewidth}\raggedright
Meaning
\end{minipage} \\
\midrule\noalign{}
\endhead
\bottomrule\noalign{}
\endlastfoot
\(S_I\) & Spin length at site \(I\) \\
\(\mathbf n_I\) & Classical spin unit direction \\
\(\mathbf S_I=S_I\mathbf n_I\) & Classical spin vector \\
\(\hat{\mathbf S}_I\) & Laboratory-frame quantum spin operator \\
\(\hat{\widetilde{\mathbf S}}_I\) & Local-frame quantum spin operator \\
\(\mathsf R_I\) & Rotation from local components to laboratory components \\
\end{longtable}

We use the rotation convention

\[
\hat{\mathbf S}_I
=\mathsf R_I\hat{\widetilde{\mathbf S}}_I,
\qquad
\mathsf R_I^{\mathsf T}\mathsf R_I=\mathsf I_3.
\]

The columns of \(\mathsf R_I\) are the local basis vectors expressed in laboratory coordinates. The exchange matrix in the local frames is

\[
\widetilde{\mathsf J}_{IJ}
=\mathsf R_I^{\mathsf T}\mathsf J_{IJ}\mathsf R_J.
\]

The polar angle \(\theta_I\) is measured from the positive \(z\) axis, and the azimuthal angle \(\phi_I\) from the positive \(x\) axis. For a general configuration, \(\mathsf R_I=\mathsf R_{i\mu}\) depends on the site. For periodic order with magnetic translation symmetry,

\[
\mathsf R_{i\mu}=\mathsf R_\mu.
\]

Accordingly, real-space expressions use \(\mathsf R_I\), while the momentum-space derivation uses the sublattice rotation \(\mathsf R_\mu\).

\subsection{Boson and Magnon Operators}

A Fourier transform changes the representation of a boson operator without changing its name. The symbol changes from \(a\) to \(b\) only when a Bogoliubov transformation replaces the sublattice basis by the magnon-band basis:

\[
\hat a_{i\mu}
\;\xrightarrow{\text{Fourier}}\;
\hat a_{\mathbf k\mu}
\;\xrightarrow{\text{Bogoliubov}}\;
\hat b_{\mathbf k n}.
\]

\begin{longtable}[]{@{}
  >{\raggedright\arraybackslash}p{(\columnwidth - 2\tabcolsep) * \real{0.5000}}
  >{\raggedright\arraybackslash}p{(\columnwidth - 2\tabcolsep) * \real{0.5000}}@{}}
\toprule\noalign{}
\begin{minipage}[b]{\linewidth}\raggedright
Symbol
\end{minipage} & \begin{minipage}[b]{\linewidth}\raggedright
Meaning
\end{minipage} \\
\midrule\noalign{}
\endhead
\bottomrule\noalign{}
\endlastfoot
\(\hat a_{i\mu}\) & Real-space Holstein--Primakoff boson \\
\(\hat a_{\mathbf k\mu}\) & Momentum-space representation of the same boson \\
\(\hat b_{\mathbf k n}\) & Positive-energy magnon annihilation operator \\
\(\hat n_{i\mu}=\hat a_{i\mu}^\dagger\hat a_{i\mu}\) & Local Holstein--Primakoff boson-number operator \\
\(n_{\mathrm B}(\varepsilon)\) & Bose--Einstein distribution function \\
\end{longtable}

The Nambu spinors before and after diagonalization are

\[
\hat\Psi_{\mathbf k}
=\begin{pmatrix}
\hat{\mathbf a}_{\mathbf k}\\
\hat{\mathbf a}_{-\mathbf k}^\dagger
\end{pmatrix},
\qquad
\hat\Phi_{\mathbf k}
=\begin{pmatrix}
\hat{\mathbf b}_{\mathbf k}\\
\hat{\mathbf b}_{-\mathbf k}^\dagger
\end{pmatrix}.
\]

Here, \(\hat{\mathbf a}_{\mathbf k}\) is a column vector in sublattice order, and \(\hat{\mathbf b}_{\mathbf k}\) is a column vector in positive-energy band order. The Fourier phase and the treatment of basis positions are specified in \lswtnoteref{LSWT derivation}{Momentum-Space BdG Hamiltonian}, subject to the convention limits listed below.

\subsection{BdG Matrices and the Nambu Metric}

\begin{longtable}[]{@{}
  >{\raggedright\arraybackslash}p{(\columnwidth - 2\tabcolsep) * \real{0.5000}}
  >{\raggedright\arraybackslash}p{(\columnwidth - 2\tabcolsep) * \real{0.5000}}@{}}
\toprule\noalign{}
\begin{minipage}[b]{\linewidth}\raggedright
Symbol
\end{minipage} & \begin{minipage}[b]{\linewidth}\raggedright
Meaning
\end{minipage} \\
\midrule\noalign{}
\endhead
\bottomrule\noalign{}
\endlastfoot
\(\mathsf H_{\mathbf k}\) & \(2N_{\mathrm{sub}}\times2N_{\mathrm{sub}}\) bosonic BdG matrix \\
\(\mathsf A_{\mathbf k}\), \(\mathsf B_{\mathbf k}\) & Normal and anomalous BdG blocks \\
\(\mathsf T_{\mathbf k}\) & Paraunitary Bogoliubov transformation matrix \\
\(\Sigma_3\) & Nambu particle--hole signature metric \\
\end{longtable}

For the particle-then-hole ordering above, the Nambu metric is

\[
\Sigma_3
\equiv
\begin{pmatrix}
\mathsf I_{N_{\mathrm{sub}}} & 0\\
0 & -\mathsf I_{N_{\mathrm{sub}}}
\end{pmatrix}
=\sigma_3\otimes\mathsf I_{N_{\mathrm{sub}}}.
\]

The subscript \(3\) refers to the Pauli matrix \(\sigma_3\), not to a matrix dimension or power. The metric extends its particle--hole grading to the \(2N_{\mathrm{sub}}\)-dimensional Nambu space. The uppercase symbol distinguishes this metric from the two-dimensional Pauli matrix and avoids using the exchange-matrix symbol \(\mathsf J\) for a second purpose. The unindexed symbol \(\Sigma\) is reserved for a possible magnon self-energy.

The metric assigns \(+1\) to the particle block and \(-1\) to the hole block, and satisfies

\[
\Sigma_3^\dagger=\Sigma_3,
\qquad
\Sigma_3^2=\mathsf I_{2N_{\mathrm{sub}}}.
\]

A paraunitary transformation preserves this metric:

\[
\mathsf T_{\mathbf k}^\dagger\Sigma_3\mathsf T_{\mathbf k}=\Sigma_3.
\]

\subsection{Energies}

\begin{longtable}[]{@{}ll@{}}
\toprule\noalign{}
Symbol & Meaning \\
\midrule\noalign{}
\endhead
\bottomrule\noalign{}
\endlastfoot
\(\varepsilon_{n\mathbf k}\) & Positive magnon energy in band \(n\) \\
\(E_{\mathrm{cl}}\) & Total classical energy \\
\(\Delta E_{\mathrm{zp}}\) & Zero-point quantum correction \\
\(E_{\mathrm{GS}}\) & LSWT ground-state energy \\
\(\mathcal E_X=E_X/N_{\mathrm{site}}\) & Energy per magnetic site \\
\end{longtable}

The total and per-site energies obey

\[
E_{\mathrm{GS}}=E_{\mathrm{cl}}+\Delta E_{\mathrm{zp}},
\qquad
\mathcal E_{\mathrm{GS}}=\mathcal E_{\mathrm{cl}}+\Delta\mathcal E_{\mathrm{zp}}.
\]

The symbol \(e_X\) is not used for energy density, to avoid confusion with the exponential base. The source symbol \(E_0\) may refer either to the total ground-state energy or to the zero-point correction; these meanings are distinguished as \(E_{\mathrm{GS}}\) and \(\Delta E_{\mathrm{zp}}\).

\subsection{Momenta and Fields}

The momentum \(\mathbf k\) labels internal magnon modes associated with the magnetic translation lattice. The momentum \(\mathbf q\) denotes an external momentum transfer, as in neutron scattering. The crystallographic and magnetic Brillouin zones are abbreviated CBZ and MBZ, respectively; the ambiguous abbreviation FBZ is not used.

The applied field \(\mathbf B_I\) is measured in tesla. The vector \(\mathbf h_I\) is the Zeeman-energy coefficient multiplying the spin operator in the Hamiltonian. The \(g\)-tensor is written as \(\mathsf g_I\). Their relation, including the explicit electron sign with positive scalar \(g\), is defined in Bilinear Spin Hamiltonian (Section~\ref{sec-lswt-bilinear-spin-hamiltonian}).

\subsection{Legacy Source Mapping}

This table translates source notation into the notation used above. The source symbols in the first column are not additional alternatives for use in the derivation.

\begin{longtable}[]{@{}
  >{\raggedright\arraybackslash}p{(\columnwidth - 4\tabcolsep) * \real{0.3333}}
  >{\raggedright\arraybackslash}p{(\columnwidth - 4\tabcolsep) * \real{0.3333}}
  >{\raggedright\arraybackslash}p{(\columnwidth - 4\tabcolsep) * \real{0.3333}}@{}}
\toprule\noalign{}
\begin{minipage}[b]{\linewidth}\raggedright
Source notation
\end{minipage} & \begin{minipage}[b]{\linewidth}\raggedright
Current notation
\end{minipage} & \begin{minipage}[b]{\linewidth}\raggedright
Reason
\end{minipage} \\
\midrule\noalign{}
\endhead
\bottomrule\noalign{}
\endlastfoot
\(i,j\) used for both cells and sites or links & Cells \(i,j\); sites \(I,J\) & Separates the index domains \\
\(\mathbf r_i\) for a cell and \(\mathbf R_I\) for a site & \(\mathbf R_i\) for a cell and \(\mathbf r_I\) for a site & Distinguishes positions from rotation matrices \\
\(J_{ij}^{\alpha\beta}\) with \(ij\) denoting a link & \(J_\ell^{\alpha\beta}\), \(\ell=(I,J)\) & Separates link identity from cell indices \\
\(\mathbf A_i\), \(A_i^{\alpha\beta}\) for single-ion anisotropy & \(\mathsf D_I\), \(D_I^{\alpha\beta}\) & Distinguishes anisotropy from the normal BdG block \(\mathsf A_{\mathbf k}\) \\
\(\mathbf R_j\) described as laboratory-to-local in the prose but used as local-to-laboratory in source equations & \(\mathsf R_I\) with \(\hat{\mathbf S}_I=\mathsf R_I\hat{\widetilde{\mathbf S}}_I\) & Follows the local-to-laboratory direction of the source equations; the conflicting source prose is recorded in the audit \\
\(\delta_{ij}\) for an undefined bond vector & \(\boldsymbol\Delta_\ell\) & Distinguishes the bond displacement from \(\boldsymbol\delta_\mu\) \\
Real-space \(a\), momentum-space \(b\), diagonal \(\beta\) & Real- and momentum-space \(a\), diagonal \(b\) & Preserves the operator name under Fourier transformation \\
\(\mathsf J\) for the Nambu metric & \(\Sigma_3\) & Distinguishes the metric from exchange matrices \\
\(\mathbf h_j\) called an external magnetic field & Applied field \(\mathbf B_I\) and Zeeman-energy vector \(\mathbf h_I\) & Separates field and energy units \\
\(R_{\mathbf k}^\alpha\) for a spin-response matrix & Spin vertex \(\mathsf V_{\mathbf k}^\alpha\) & Distinguishes the vertex from rotations; its precise definition and dagger convention remain unspecified \\
\(E_{\mathbf k,\mu}\), \(\varepsilon_{n,\mathbf k}\) & \(\varepsilon_{n\mathbf k}\) & Distinguishes band and sublattice indices \\
\(E_0\) & \(E_{\mathrm{GS}}\) or \(\Delta E_{\mathrm{zp}}\) & Separates the two source meanings \\
\(L\), \(N\), \(m_s\) & \(N_{\mathrm{uc}}\), \(N_{\mathrm{site}}\), \(N_{\mathrm{sub}}\) & Fixes the meaning of each system size \\
\end{longtable}

\subsection{Conventions Not Yet Fixed}

The following quantities or conventions are not fully specified. Reserving a symbol does not establish the associated physical definition.

\begin{itemize}
\tightlist
\item
  Fourier-transform sign and full-position versus periodic gauge.
\item
  Folding and normalization between the CBZ and MBZ.
\item
  Momentum-dependent magnon-band ordering and band identity.
\item
  Symbols and dimensions for the positive-band energy matrix and the full Nambu diagonal matrix.
\item
  Exact definitions and dagger conventions for the spin-response matrices \(\mathsf V_{\mathbf k}^\alpha\), \(\mathsf U^\alpha\), and \(\mathsf S_{\mathbf k}\).
\item
  Phase and contraction conventions for local circular components and their response formulas.
\item
  Normalization of zero-point momentum sums and normal-ordering constants.
\item
  Units for length, real time, and thermal Hall conductivity.
\end{itemize}

\subsection{References}

\subsubsection{Internal Documents}

\begin{itemize}
\tightlist
\item
  Bilinear Spin Hamiltonian (Section~\ref{sec-lswt-bilinear-spin-hamiltonian}): defines the exchange, anisotropy, Zeeman, and link-counting equations that use this notation.
\item
  Classical Order and Local Frame (Section~\ref{sec-lswt-classical-order-and-local-frame}): defines the rotation and local-frame conventions.
\item
  \lswtnoteref{LSWT derivation}{Momentum-Space BdG Hamiltonian}: defines the Fourier, Nambu-spinor, and BdG-block definitions.
\item
  \lswtnoteref{LSWT derivation}{Paraunitary Diagonalization}: defines the operational use of the Nambu metric and paraunitary transformation.
\item
  \lswtnoteref{LSWT observables}{Thermodynamics}: defines the thermodynamic definitions that use the reserved energy and entropy symbols.
\end{itemize}

\subsubsection{External Sources}

\begin{itemize}
\tightlist
\item
  None.
\end{itemize}
